\documentclass[11pt,a4paper]{article}
\usepackage{amsmath,amssymb,amsthm,mathtools,microtype}
\usepackage[margin=26mm]{geometry}
\usepackage[colorlinks=true,linkcolor=blue,citecolor=blue,urlcolor=blue]{hyperref}
\newtheorem{theorem}{Theorem}[section]
\newtheorem{lemma}[theorem]{Lemma}
\newtheorem{corollary}[theorem]{Corollary}
\newtheorem{proposition}[theorem]{Proposition}
\newcommand{\Q}{\mathbb Q}\newcommand{\Z}{\mathbb Z}
\DeclareMathOperator{\Tr}{Tr}\DeclareMathOperator{\Num}{Num}
\setlength{\emergencystretch}{2em}
\title{Two height tests and a three-square interval\protect\\
in the Bhargava--Burungale construction}
\author{ZQ Diophantine Project}
\date{23 November 2026}
\begin{document}
\maketitle
\begin{abstract}
We transfer the project's aggregate-trace height method to the carrier
over $\Q(\sqrt{13})$ in Bhargava--Burungale, arXiv:2611.16130v2.
This reduces their five arithmetic tests for squaring to two, with the
same local predicate. A fixed power of one rational trace also removes
the need for their large elliptic multiplication map. Independently,
a fixed $2$-adic coset permits
three rational squares in the final interval condition in place of
four. Relative to the same quadratic presentation of their set $T$,
the quartic wrapper therefore uses $r+18$, rather than $r+19$,
auxiliary rational variables. The parameter $r$ is not evaluated;
these results do not assert a new numerical defining-variable record.
All assertions retain the stated Bhargava--Burungale arithmetic inputs.
The local construction also defines every fixed quadratic relation
directly; in particular multiplication requires two tests, without
three calls to a squaring relation.
\end{abstract}

\section{Data and dependence}
Put
\[
 F=\Q(\sqrt{13}),\qquad d=\frac{1+\sqrt{13}}2,\qquad
 E/F:y^2=x^3-d^2x,\qquad \Gamma=2E(F).
\]
Write $\sigma$ for the nontrivial automorphism of $F/\Q$.
Bhargava--Burungale \cite[Lemmas 4.1--4.2]{BB} prove that
$\Gamma$ is infinite cyclic and Diophantine. Fix its generator $P_0$.
For $P_j=n_jP_0$, put $\Delta=n_1^2-n_2$.

Use the local predicate $\mathcal C(P_1,P_2;\xi)$ of
\cite[Definition 5.3]{BB} without alteration. Let
\[
 G(X)=(X^2-1)(X^2-9)(X^2-25),\qquad
 M(\xi)=\prod_{\substack{p>25\\v_p(G(\xi))>0\text{ odd}}}p.
\]
The two imported conclusions required below are
\begin{align}
 h(\xi)&\le H M(\xi)^{24}\quad(G(\xi)\ne0),\label{eq:height}\tag{H}\\
 \mathcal C(P_1,P_2;\xi)&\Longrightarrow M(\xi)^{24}\mid\Delta.
 \label{eq:div}\tag{D}
\end{align}
They are \cite[Proposition 3.1 and Lemma 5.4]{BB}. The construction
of witnesses when $\Delta=0$ is their Lemma 5.2 and the argument of
Lemma 5.6. We explicitly check below that the latter applies to any
fixed finite collection of rational shifts.

The additional published height input is von Buhren's family theorem
\cite[Theorem 1.1, p.~308]{VB}. The argument in the next section is
the specialization to this carrier of the project's positive
weighted-trace argument \cite{Weighted}. Thus the aggregate-trace
method is prior project work; the present result is its transfer to
the different carrier and its use to shorten the BB test list.

\section{Aggregate trace on the BB carrier}
Define $x^*(O)=0$ and $x^*(P)=x(P)$ for $P\ne O$. This is a
piecewise definition, not the value of the rational function $x$ at
its pole. Its use in a Diophantine relation means a finite disjunction
according to which points are $O$.

\begin{theorem}\label{thm:trace}
There are constants $A>0$, $J\ge0$ such that for all integers
$k_1,k_2$,
\[
 k_1^2+k_2^2\le A h\!\left(
 \Tr_{F/\Q}\bigl(x^*(k_1P_0)+x^*(k_2P_0)\bigr)
 \right)+J.
\]
The constants are independent of the indices. The additive constant
is not made effective by the proof.
\end{theorem}
\begin{proof}
At the embedding $\sqrt{13}>0$, set
$d_+=d$ and $d_-=-d^\sigma$. Both are positive, and $d_+d_-=3$.
Their squareclasses are independent in $F(i)^\times/F(i)^{\times2}$.
Indeed, an element of the real field $F$ that is a square in $F(i)$
is either a square or minus a square in $F$. Neither $d_+$ nor
$-d_+$ is a square in $F$, by their real signs; the same holds for
$d_-$. Their product is $3$, and neither $3$ nor $-3$ is a square in
$F$ (the former would give $\Q(\sqrt3)\subseteq F$).

Over $F(i,\sqrt{d_+},\sqrt{d_-})$, identify both curves with
$E_0:Y^2=X^3-X$ by division of the coordinates by
$(d_+,d_+^{3/2})$ and $(d_-,d_-^{3/2})$. Let $R,R'$ be the
images of $P_0,P_0^\sigma$. They are independent modulo torsion
over $\Z[i]$. An automorphism fixing $F(i,\sqrt{d_+})$ and
negating $\sqrt{d_-}$ fixes $R$ and negates $R'$; subtracting the
conjugate of a proposed relation proves independence, since a nonzero
endomorphism of an elliptic curve is an isogeny.

First omit the indices that are zero. There are only finitely many
such patterns; suppose $k\in\{1,2\}$ indices remain. On $E_0^{2k}$
consider
\[
 \mathcal F=d_+\sum_{j=1}^k x_j+d_-\sum_{j=1}^k x'_j.
\]
At our point the coordinates are $[k_j]R,[k_j]R'$, and
$\mathcal F$ is the required rational trace. No positive-dimensional
abelian translate through this point, meeting the domain of
$\mathcal F$, can have constant $\mathcal F$. To prove this, take an
elliptic direction in such an abelian subvariety and parametrize it
by $E_0$ with coordinate endomorphisms $\alpha_j,\alpha'_j\in\Z[i]$.
If two first-block coordinates have a common pole, eliminating the
parameter gives
\[
 (k_i\alpha_j-k_j\alpha_i)R=O,
 \qquad\text{hence}\qquad k_i\alpha_j=k_j\alpha_i.
\]
The leading coefficient of their combined double pole, including all
first-block coordinates with that pole, is therefore
\[
 \frac{d_+}{\alpha_i^2}\sum_j(k_i/k_j)^2\ne0.
\]
The sum is a positive rational number. A common pole between the two
blocks would give a nontrivial $\Z[i]$-relation between $R$ and $R'$,
which is impossible. The same positive-sum argument applies within
the second block. At least one coordinate of the direction is
nonconstant, so there is an uncancelled pole.

Take the integral reduced closure of the graph of $\mathcal F$ in
$E_0^{2k}\times\mathbf P^1$. A positive-dimensional abelian translate
contained in a fiber and passing through our point would meet the
rational-function domain densely, contrary to the preceding paragraph.
The points under consideration belong to the fixed finite-rank
subgroup generated by $R$ in the first block and $R'$ in the second.
They belong to the set $C(\Gamma',\epsilon)$ of \cite{VB} with error
point zero. Thus the subgroup does not vary with the indices.

For completeness, embed each $E_0$ by $\mathcal O(3O)$ and its
product by Segre. Restriction of sections in every positive degree is
surjective: for one factor this follows from
\[
0\to\mathcal O_{\mathbf P^2}(n-3)\to\mathcal O_{\mathbf P^2}(n)
 \to\mathcal O_{E_0}(3nO)\to0
\]
and $H^1(\mathbf P^2,\mathcal O(n-3))=0$; take tensor products for
the product embedding. This checks condition (1.1) of \cite{VB}.
The line bundle is ample and symmetric, and the parameter space is
the projective line with its standard height. The theorem bounds the
canonical product height by $A_0h(\mathcal F)+J_0$.
Pointed-isomorphism invariance and Galois invariance identify that
height with a fixed positive multiple of
$2\sum_j k_j^2\widehat h(P_0)$. Divide by this positive coefficient
and take the maximum constants over the zero-index patterns. The
all-zero pattern is immediate.
\end{proof}

\section{Prime progressions for arbitrary rational shifts}
\begin{lemma}\label{lem:progression}
Let $\mathcal P$ be a set of primes of positive relative upper
density, let $B,\Theta\subset\Q$ be finite, and let $L_0>0$.
There exist $s\in\Q$ and a positive integer $T$ such that
\[
 T(s+\beta-\theta)\in\mathcal P,\qquad
 s+\beta-\theta>L_0
 \quad(\beta\in B,\theta\in\Theta),
\]
and every displayed prime is coprime to $T$. For each fixed $\beta$,
distinct $\theta$ give distinct primes. One may additionally require
that $T$ be divisible by any prescribed positive integer.
\end{lemma}
\begin{proof}
Choose a positive integer $D$ divisible by the prescribed integer and
clearing all denominators in $B-\Theta$. Choose an integer $m$ so that
\[
 r_{\beta,\theta}=m+D(\beta-\theta)>
 \max(DL_0,1)
\]
for every pair. Discard the finitely many primes at most $D$ from
$\mathcal P$. By Green--Tao \cite[Theorem 1.2]{GT}, this set contains
an arithmetic progression $q_0,q_0+v,\ldots,q_0+Nv$ with $v>0$ and
$N\ge\max r_{\beta,\theta}$. Put
\[
 T=Dv,\qquad s=\frac{q_0+mv}{Dv}.
\]
Then $T(s+\beta-\theta)=q_0+r_{\beta,\theta}v=p$ is a prime
in $\mathcal P$. We have $p>D$ and $p>v$, so $p$ divides neither
$D$ nor $v$. This is the needed denominator-cancellation check.
The lower bound and distinctness follow directly.
\end{proof}

\begin{proposition}\label{prop:finite}
If $\Delta=0$, then for every finite $B\subset\Q$ there is
$s\in\Q$ such that $\mathcal C(P_1,P_2;s+\beta)$ holds for all
$\beta\in B$, with separate witnesses at the different shifts.
\end{proposition}
\begin{proof}
Use Lemma~\ref{lem:progression} with primes congruent to $17$ modulo
$312$ and exceeding $25$, with $\Theta=\{-5,-3,-1,1,3,5\}$,
$L_0=2$, and $14\mid D$. At each shift, the sextic is the product
of six distinct chosen primes divided by $T^6$, with no cancellation.
Consequently the verification in \cite[Lemma 5.6]{BB} applies at
each shift separately. Its six triples handle those six primes; at
$2,7$ the denominator has exactly the same preparation as in that
proof. Nothing in that verification uses the special origin or the
number of the shifts.
\end{proof}

\section{Two tests for squaring}
For a positive integer $m$, put
\[
 T_m=\Tr_{F/\Q}\bigl(x^*(mP_1)+x^*(mP_2)\bigr).
\]
Theorem~\ref{thm:trace} gives
\begin{equation}\label{eq:amp}
 m^2(n_1^2+n_2^2)\le A h(T_m)+J,
\end{equation}
with the same $A,J$ for every $m$.

\begin{theorem}\label{thm:two}
Fix an integer $m\ge1$ such that
\[
 m^2>2AH+A\log 2+J.
\]
Then $n_2=n_1^2$ is equivalent to the existence of $s\in\Q$ with
\[
 \mathcal C(P_1,P_2;s),\qquad
 \mathcal C(P_1,P_2;s+T_m).
\]
\end{theorem}
\begin{proof}
If $\Delta\ne0$, \eqref{eq:height} and \eqref{eq:div} bound each
of $h(s),h(s+T_m)$ by $H|\Delta|$. Subtracting gives
\[
 h(T_m)\le2H|\Delta|+\log2.
\]
Let $S=n_1^2+n_2^2\ge1$. Since the indices are integers,
$|\Delta|\le n_1^2+|n_2|\le S$. Substitution in
\eqref{eq:amp} gives
$m^2S\le(2AH+A\log2+J)S$, a contradiction. Conversely,
Proposition~\ref{prop:finite} applies to $B=\{0,T_m\}$.
\end{proof}

The identity case may instead be handled by the disjunction
$(P_1,P_2)=(O,O)$, or $P_1,P_2\ne O$ and the two displayed tests.
This avoids any need to evaluate the abscissa at $O$ in the nonzero
branch. All the operations are Diophantine because $m$ is fixed.

\subsection{Power amplification without elliptic multiplication}
The positivity of doubled points permits a further simplification.
Put
\[
 T=\Tr_{F/\Q}\bigl(x^*(P_1)+x^*(P_2)\bigr),\qquad
 A_*=A+\frac{J}{\log3}.
\]
If $(P_1,P_2)\ne(O,O)$, then $T>\sqrt{13}>3$. Indeed, at
each real embedding a nonzero doubled point has abscissa greater
than the absolute value of the corresponding conjugate of $d$;
the trace of each nonzero summand is greater than
$d-d^\sigma=\sqrt{13}$. Hence $h(T)\ge\log3$, and
Theorem~\ref{thm:trace} yields
\begin{equation}\label{eq:noerror}
 n_1^2+n_2^2\le A_* h(T).
\end{equation}

\begin{theorem}\label{thm:power}
Choose any fixed positive integer $k$ satisfying
$k>A_*(2H+\log2)$. The squaring relation is equivalent to
\[
 \exists s\in\Q:\quad
 \mathcal C(P_1,P_2;s),\qquad
 \mathcal C(P_1,P_2;s+T^k).
\]
One may choose $k$ to be a power of two. No large elliptic
multiplication map occurs in this formula.
\end{theorem}
\begin{proof}
If $\Delta\ne0$, set $S=n_1^2+n_2^2\ge1$ and combine
\eqref{eq:height}, \eqref{eq:div}, and \eqref{eq:noerror}:
\[
 \frac{k}{A_*}S\le k h(T)=h(T^k)
 \le 2H|\Delta|+\log2\le(2H+\log2)S,
\]
contrary to the choice of $k$. If $\Delta=0$, use
Proposition~\ref{prop:finite} with $B=\{0,T^k\}$, including
$T=0$ in the all-zero case.
\end{proof}

For an unrestricted-degree formula, substituting $T^k$ introduces
no new retained rational parameter. For a quadratic circuit,
choosing $k=2^j$ computes this power by $j$ successive squarings.
This is a circuit prescription rather than a numerical bound on $j$:
the constants have not been evaluated.

The original proof has five copies of $\mathcal C$, each containing
six Gaussian point triples. The new proof has two copies: the number
of triples drops from $30$ to $12$, and the number of auxiliary
Gaussian points $U_j$ from $90$ to $36$. These are counts of blocks
and points, not counts of the final rational coordinates after
encoding group membership, division, and DPRM arithmetic.

\section{Direct quadratic relations and multiplication}
The local argument is not specific to squaring. Fix
$F_0\in\Z[X_1,\ldots,X_s]$ of total degree at most two. For
$P_j=n_jP_0$ ($1\le j\le s$), put
$\Delta_{F_0}=F_0(n_1,\ldots,n_s)$ and keep $n_0=1$.
We describe a predicate $\mathcal C_{F_0}(P_1,\ldots,P_s;\xi)$.
It is the BB predicate with the following explicit changes only.

For each $1\le\ell\le6$, choose $U_j\in\Gamma+[i]\Gamma$
for $0\le j\le s$, put $V_j=P_j-U_j$, and form the formal-log
ratios, separately for each tuple,
\[
 B_j=\frac{f(z(U_j))}{f(z(U_0))},\qquad
 C_j=\left(\frac{f(z(V_j))}{f(z(V_0))}\right)^\sigma,
 \qquad 1\le j\le s.
\]
Here $z=-x/y$, $\rho=1/x$, their values at $O$ are zero, and
$f$ is the degree-$21$ truncation in BB (5.2). Require the reference
ratios to be defined. With $G(\xi)\ne0$, $R=14G(\xi)$, and
$H_R(X)=R/(R+X^2)$, use
\begin{align*}
 W_\ell={}&
 \prod_{j=0}^{s}H_R(\rho(U_j))H_R(\rho(V_j^\sigma))\,
 \prod_{j=1}^{s}H_R\!\left(\frac{B_j-C_j}{R^{23}}\right)\\
 &\hspace{20mm}\cdot
 H_R\!\left(\frac{F_0(B_1,\ldots,B_s)}{R^{23}}\right).
\end{align*}
There are $3(s+1)$ factors. Require rational $c,a_\xi,b_\xi,z_\xi$
such that
\begin{gather*}
 w_0=1-\operatorname{N}_{F(i)/\Q}
      \left(\prod_{\ell=1}^{6}(1-W_\ell)\right),\qquad
 c(R+w_0^2)^2=2R,\\
 (1-c)(a_\xi^4-c^2)(b_\xi^4-c^2)=z_\xi^2.
\end{gather*}
Every quotient in these expressions is required to be defined.

\begin{lemma}\label{lem:quadraticlocal}
The predicate $\mathcal C_{F_0}$ implies
$M(\xi)^{24}\mid\Delta_{F_0}$. If $\Delta_{F_0}=0$, then for
every finite set $B\subset\Q$ there is a common shift $v\in\Q$
such that $\mathcal C_{F_0}(P_1,\ldots,P_s;v+\beta)$ holds for
every $\beta\in B$.
\end{lemma}
\begin{proof}
Fix $p>25$ and a place $\mathfrak p$ of $F(i)$ above $p$, with
$v(p)=1$. The required local conditions are
\begin{gather*}
 v(\rho(U_j)),\ v(\rho(V_j^\sigma))\ge1
       \quad(0\le j\le s),\\
 B_j-C_j\equiv0\pmod{\mathfrak p^{24}}
       \quad(1\le j\le s),\qquad
 F_0(B_1,\ldots,B_s)\equiv0\pmod{\mathfrak p^{24}}.
\end{gather*}
Write $U_j=[\gamma_j]P_0$, $V_j=[\delta_j]P_0$ with Gaussian
integers $\gamma_j+\delta_j=n_j$. In particular,
$\gamma_0+\delta_0=1$. The first local conditions put all points
in the appropriate formal groups. BB's relative logarithm estimate
$f(z(Q))/\log_E(z(Q))\in1+\mathfrak p^{24}$ gives
\[
 \gamma_0B_j\equiv\gamma_j,\qquad
 \delta_0C_j\equiv\delta_j\pmod{\mathfrak p^{24}}.
\]
The errors here are integral Gaussian indices times elements of
$\mathfrak p^{24}$, so this assertion does not assume that $B_j$
or $C_j$ is integral. Using the comparison condition gives
\[
 B_j=(\gamma_0+\delta_0)B_j
      \equiv\gamma_j+\delta_j=n_j\pmod{\mathfrak p^{24}}.
\]
Thus $B_j$ is integral. Evaluation of the integral polynomial $F_0$
now gives $\Delta_{F_0}\equiv0\pmod{p^{24}}$. Zero numerator
points cause no exception: their index and logarithm ratio are both
zero, while the two reference denominators remain guarded.

Suppose $p\mid M(\xi)$ and let $w=v_p(R)>0$, which is odd.
At every place above $p$, all $H_R$-values are integral because
\[
 v(H_R(X))=\max(0,w-2v(X)).
\]
If a formal-group condition fails, its $\rho$-value has valuation
at most zero. If a congruence fails, its normalized argument has
valuation at most $23-23w\le0$. Hence a failure makes the relevant
$W_\ell$ divisible by $\mathfrak p^w$. If every tuple failed at
every place above $p$, the product inside the norm would be $1$
modulo $p^w$ in the integral closure, so its norm would be $1$
modulo $p^w$. This would give $v_p(w_0)\ge w$, then
$v_p(R+w_0^2)=w$ and $v_p(c)=-w$, contradicting the square-denominator
implication of BB Lemma 2.2. Some tuple succeeds above each such
$p$, proving $M(\xi)^{24}\mid\Delta_{F_0}$.

For the converse, suppose $\Delta_{F_0}=0$ and select the common
shift by Lemma~\ref{lem:progression}, exactly as in
Proposition~\ref{prop:finite}. Fix one of the six primes at one
shift. It is congruent to $17$ modulo $312$, so it splits completely
in $F(i)$, the reduction is ordinary, and the conjugate curve is a
nonsquare quadratic twist. At a selected place its two Frobenius
endomorphisms are $[\pi]$ and $[-\pi]$, with
$\pi\in\Z[i]$, $\pi\overline\pi=p$, and
$\pi\equiv1\pmod{2\Z[i]}$. Put
\[
 U_j=\left[\frac{1-\pi}{2}\right]P_j,
 \qquad V_j=\left[\frac{1+\pi}{2}\right]P_j
 \quad(0\le j\le s).
\]
Since $P_j=[2]Q_j$ with $Q_j\in E(F)$, these become
$[1-\pi]Q_j$ and, after conjugation, $[1+\pi]Q_j^\sigma$.
Both reduce to $O$. The exact logarithm ratios are $n_j$, so their
truncations are congruent to $n_j$ modulo $p^{24}$. Integral
polynomial evaluation proves the last local congruence. The
multipliers are nonzero, and the formal logarithm is injective here,
so the reference denominators are nonzero.

At this chosen prime $v_p(R)=1$; every argument in the displayed
$W_\ell$ has positive valuation. Thus $W_\ell\equiv1$ at the
chosen place. All factors for every tuple are integral above $p$,
so the norm of their product of complements is divisible by $p$.
Consequently $w_0\equiv1\pmod p$ and $v_p(c)=1$. There are still
only six selected primes at each shift, irrespective of $s$.

Finally, $R=14p_1\cdots p_6/T^6>2$, with $14\mid T$. At $2,7$,
write $w=v_\ell(R)\le-5$, odd. Then
$v_\ell(R+w_0^2)=\min(w,2v_\ell(w_0))\le w$, so $c$ has
positive valuation there, and $v_2(c)\ge6$. Its valuations outside
$2,7,p_1,\ldots,p_6$ are even, its squareclass is
$7p_1\cdots p_6\equiv7\pmod8$, and $0<c\le2/R<1$.
The odd valuation of $R$ at a split selected prime prevents
$R+X^2$ from vanishing for $X\in F(i)$; the rational denominator
$R+w_0^2$ is also positive. BB Lemma 2.2 supplies
$a_\xi,b_\xi,z_\xi$. This proves all definedness, valuation and
real requirements at every shift.
\end{proof}

The proof of Theorem~\ref{thm:trace} works for every fixed finite
number of indices, with constants depending on that number: the same
pole argument excludes all positive-dimensional abelian translates.
Apply it also to the fixed extra point $P_0$ and set
\[
 T=\Tr_{F/\Q}\left(x(P_0)+\sum_{j=1}^{s}x^*(P_j)\right),
 \qquad S=1+\sum_{j=1}^{s}n_j^2.
\]
Then $T>\sqrt{13}>3$, and there is $A_{*,s}>0$ with
$S\le A_{*,s}h(T)$. Let $C_{F_0}$ be the sum of the absolute
values of the coefficients of $F_0$, increased to at least one.
Every monomial of degree at most two has absolute value at most $S$,
so $|\Delta_{F_0}|\le C_{F_0}S$. Any smaller positive constant
with this proved bound may be used instead.

\begin{theorem}\label{thm:quadratic}
Choose a fixed positive integer
$k>A_{*,s}(2HC_{F_0}+\log2)$. Then
$F_0(n_1,\ldots,n_s)=0$ is equivalent to
\[
 \exists v\in\Q:\quad
 \mathcal C_{F_0}(P_1,\ldots,P_s;v),\qquad
 \mathcal C_{F_0}(P_1,\ldots,P_s;v+T^k).
\]
\end{theorem}
\begin{proof}
If $\Delta_{F_0}\ne0$, the two tests bound $h(T^k)$ by
$2H|\Delta_{F_0}|+\log2$. This contradicts
$h(T^k)=kh(T)\ge kS/A_{*,s}$ and the choice of $k$.
The converse is Lemma~\ref{lem:quadraticlocal} with shifts
$0,T^k$.
\end{proof}

In particular, use $F_0=X_1X_2-X_3$ to define multiplication
directly. Here the sharper constant $C_{F_0}=1/2$ is valid, since
\[
 2|n_1n_2-n_3|\le n_1^2+n_2^2+n_3^2+1=S.
\]
Thus $k>A_{*,3}(H+\log2)$ suffices. It uses two predicate blocks
with six Gaussian quadruples
each, hence $48$ auxiliary Gaussian points. Three calls to the
two-test squaring relation use $108$ such points, while the original
five-test implementation uses $270$. These figures count only the
$U_j$ points and do not include point encodings, intermediate outputs,
or the final rational variables. Addition remains the elliptic group
law. The direct quadratic interface is useful when compiling the
DPRM arithmetic, but its full coordinate cost still requires a
separate compiler calculation.

\section{Three squares for the final interval}
Write
\[
 y(P)=a(P)+b(P)\sqrt{13}\qquad(P\in\Gamma\setminus\{O\}).
\]
\begin{lemma}\label{lem:coset}
There is a coset $\mathcal K=P_*+[N]\Gamma$ such that
\[
 v_2(a(P))=-2,\qquad v_2(b(P))=-3\qquad(P\in\mathcal K),
\]
and the pairs $(a(P),b(P))$, $P\in\mathcal K$, are dense in
$\mathbb R^2$.
\end{lemma}
\begin{proof}
Take $Q=(d+1,d+2)\in E(F)$ and $P_*=2Q$. Direct duplication gives
\[
 x(P_*)=\frac{187+19\sqrt{13}}{72},\qquad
 y(P_*)=-\frac{205}{108}-\frac{191}{216}\sqrt{13}.
\]
Thus the two valuations at $P_*$ are $-2,-3$. They persist in a
sufficiently small neighborhood of $P_*$ in $E(F\otimes\Q_2)$.
Choose an open subgroup $U$ of the compact local group with
$P_*+U$ in this neighborhood. The quotient by $U$ is finite, so a
positive integer $N$ has $[N]\Gamma\subseteq U$.

The independence modulo torsion used in \cite[Lemma 4.2]{BB}
persists after multiplication by $N$. Consequently the real closure
of $[N]\Gamma$ is still the product of the two identity circles.
Translating by $P_*\in\Gamma$ leaves that closure unchanged.
The same coordinate argument as in that lemma gives the asserted
density in $\mathbb R^2$.
\end{proof}

\begin{lemma}\label{lem:three}
For every $q\in\Z$ and every $\epsilon>0$, there is
$P'\in\Gamma\setminus\{O\}$ such that
\[
 a(P')<q<b(P'),\qquad b(P')-a(P')<\epsilon,
\]
and $(q-a(P'))(b(P')-q)$ is a sum of three rational squares.
\end{lemma}
\begin{proof}
Choose $P'$ from the dense coset of Lemma~\ref{lem:coset} with
the displayed real inequalities. Since $v_2(q)\ge0$, the positive
rational number $r=(q-a(P'))(b(P')-q)$ has $v_2(r)=-5$.
Write $r=A/B$ in lowest terms, with $A,B>0$. The integer $AB$ has
odd $2$-adic valuation, so is not of the form $4^a(8b+7)$.
The Legendre--Gauss three-square theorem, in the form recorded in
\cite{Squares}, writes $AB=x_1^2+x_2^2+x_3^2$ with integers $x_i$.
Then $r=\sum_i(x_i/B)^2$.
\end{proof}

Fix $C_0$ as in \cite[Theorem A]{BB}, and keep exactly its original
set $T_{C_0}$. In particular, no extra coset condition is added to
the definition of $T$ or to its quadratic presentation.

\begin{theorem}\label{thm:wrapper}
In \cite[Theorem A]{BB}, the condition $a\le q\le b$ can be
replaced by
\[
 (q-a)(b-q)=v_1^2+v_2^2+v_3^2,\qquad v_i\in\Q.
\]
If $T_{C_0}$ has a quadratic presentation with $r$ auxiliary
variables, the resulting definition of $\Z$ has a quartic
presentation with at most $r+18$ auxiliary variables besides $q$.
\end{theorem}
\begin{proof}
Since $(u,a,b)\in T$ implies $a\le b$, the new condition implies
$a\le q\le b$. Thus the original soundness proof applies unchanged.
For completeness, choose the point supplying $u$ exactly as in BB,
then apply Lemma~\ref{lem:three} with
$\epsilon=C_0^{-\Num(u)^{24}}$. The interval still contains the
integer $q$, hence defines an element of the same $T$.

Keep the nine intermediates in \cite[Remark 6.1]{BB}, the six
variables $u,a,b,s_0,t_0,z_0$, and the $r$ variables in the chosen
quadratic presentation of $T$. Add the three $v_i$ instead of four.
Every equation is quadratic, including the new interval equation.
Their sum of squares is quartic and the count is $r+9+6+3=r+18$.
\end{proof}

\section{Scope and further work}
Neither shortening above supplies a rational point on the project's
one-parameter actual detector fibers. In particular it does not
identify the genus-one quotient with its required physical lift.
The integers $k,C_0$ and the DPRM presentation are fixed existential
data here; no numerical value of $r$ is claimed. The two-test
construction can reduce the cost of compiling that arithmetic model,
but an actual compiler and proof-preserving interface comparison are
needed before advertising any final variable count.

The companion checker verifies the duplication seed, its valuation
identity for every integer input, the split-sextic identity, the
formal-log coefficients used by BB, and sample congruence-class
identities by exact arithmetic. It does not prove the cited height,
density, three-square, or prime progression theorems.

\begin{thebibliography}{9}
\bibitem{BB} M.~Bhargava and A.~Burungale,
\emph{A Diophantine definition of $\Z$ in $\Q$},
arXiv:2611.16130v2, 19 November 2026.
\url{https://arxiv.org/html/2611.16130v2}.
\bibitem{VB} J.~von Buhren,
\emph{Borne de hauteur semi-effective pour le probl\`eme de
Mordell--Lang dans une vari\'et\'e ab\'elienne},
J.~Th\'eor. Nombres Bordeaux \textbf{29} (2017), 307--320,
Theorem 1.1.
\url{https://www.numdam.org/article/JTNB_2017__29_1_307_0.pdf}.
\bibitem{Weighted} ZQ Diophantine Project,
\emph{Heights of positive weighted elliptic traces}, 14 November 2026,
\nolinkurl{papers/trace_heights/weighted_trace_20261114_v27.tex}.
\bibitem{GT} B.~Green and T.~Tao,
\emph{The primes contain arbitrarily long arithmetic progressions},
Ann. of Math. \textbf{167} (2008), 481--547, Theorem 1.2.
\url{https://annals.math.princeton.edu/wp-content/uploads/annals-v167-n2-p03.pdf}.
\bibitem{Squares} L.~Goldmakher and P.~Pollack,
\emph{Refinements of Lagrange's four-square theorem},
statement of the Legendre--Gauss three-square theorem.
\url{https://www.pollack-math.net/4squares-arxiv.pdf}.
\end{thebibliography}
\end{document}
